% SEGMENT OLP-0687-S01
% Part: proof-theory
% Chapter: propositions-as-types
% Section: normalization

\documentclass[../../../include/open-logic-section]{subfiles}

\begin{document}

\olfileid{int}{pty}{nor}

\olsection{இயல்பாக்கம்}

இப்பிரிவில், பொருத்தமான ஒரு குறைத்தல் வரிசையைத்
தேர்ந்தெடுத்தால் எந்த வருவித்தலையும் இயல்பு வடிவ
வருவித்தலாகக் குறைக்கலாம் என்பதை நிறுவுகிறோம்.
இதனை \emph{இயல்பாக்கப் பண்பு} என்பார்கள்.
கூற்றுகள் வகைகளாகும் பொருத்தத்தைப் பயன்படுத்துவோம்:
ஒவ்வொரு நிறுவல் சொல்லையும் இயல்புவடிவமாகக்
குறைக்கலாம் என்று காட்டினால், இயல்பான வருவித்தலின்
!!{derivation}s க்கும் இயல்பாக்கம் பின்வரும்.

முதலில், சொற்களின் சிக்கலை அளக்கும் சில சார்புகளை
வரையறுப்போம். !!a{formula} இன் \emph{நீளம்}
~$\len{!A}$ பின்வருமாறு வரையறுக்கப்படுகிறது:
\begin{align*}
  \len{p} &= 0\\
  \len{!A \land !B} &= \len{!A} + \len{!B} + 1\\
  \len{!A \lor !B} &= \len{!A} + \len{!B} + 1 \\
  \len{!A \lif !B} &= \len{!A} + \len{!B} + 1.
\end{align*}
ரெடெக்ஸ்~$M$ இன் சிக்கலை அதன் \emph{வெட்டு
நிலை}~$\cutrank{M}$ கொண்டு அளக்கிறோம்:
\begin{align*}
  \cutrank{(\lambd[\typeof{x}{!A}][\typeof{N}{!B}])Q} &=
  \len{!A} + \len{!B} + 1 \\
  \cutrank{\ande{i}{\andi{\typeof{M}{!A}}{\typeof{N}{!B}}}} &=
  \len{!A} + \len{!B} + 1 \\
  \cutrank{\ore{\ori{i}{!A_{i}}{\typeof{M}{!A_i}}}
    {\typeof{x_1}{!A_1}}{\typeof{N_1}{!C}}
    {\typeof{x_2}{!A_2}}{\typeof{N_2}{!C}}} &=
  \len{!A_1} + \len{!A_2} + 1
\end{align*}
ஒரு நிறுவல் சொல்லின் சிக்கல், அதிலுள்ள ரெடெக்ஸ்களின்
வெட்டு நிலைகளில் மிகப் பெரியது; அது இயல்புவடிவமாக
இருந்தால் சிக்கல் $0$. எனவே, கீழுள்ள தொகுதி வெறுமையாக
இருக்கும்போது அதன் மிகப் பெரிய மதிப்பைப் பூச்சியம் என்று
கொள்ளும் மரபுடன்,
\[
  \maxrank{M} = \max \{\cutrank{N} | N \text{ என்பது } M \text{ இன் உட்சொல்;
    அது ரெடெக்ஸ்}\}
\]
என்று எழுதுகிறோம்.

% SEGMENT OLP-0687-S02
\begin{lem}\ollabel{lem:subst}
  $\Subst{M}{\typeof{N}{!A}}{\typeof{x}{!A}}$ ஒரு
  ரெடெக்ஸாகவும் $M \not\ident x$ ஆகவும் இருந்தால்,
  பின்வரும் நேர்வுகளில் ஒன்று அமையும்:
  \begin{enumerate}
  \item $M$ தானே ஒரு ரெடெக்ஸ்; அல்லது
  \item $M$ ஆனது $\ande{i}{x}$ என்ற வடிவிலும்,
    $N$ ஆனது $\andi{P_1}{P_2}$ என்ற வடிவிலும் இருக்கும்;
  \item $M$ ஆனது $\ore{x}{x_1}{P_1}{x_2}{P_2}$ என்ற
    வடிவிலும், $N$ ஆனது $\ori{i}{}{Q}$ என்ற வடிவிலும் இருக்கும்;
  \item $M$ ஆனது $x Q$ என்ற வடிவிலும், $N$ ஆனது
    $\lambd[x][P]$ என்ற வடிவிலும் இருக்கும்.
  \end{enumerate}
  முதல் நேர்வில் $\cutrank{\Subst{M}{N}{x}} = \cutrank{M}$;
  மற்ற நேர்வுகளில் $\cutrank{\Subst{M}{N}{x}} = \len{!A}$.
\end{lem}

% SEGMENT OLP-0687-S03
\begin{proof}
  $M$ இன் கட்டமைப்பின்மீது தொகுத்தறிதலால் நிறுவுகிறோம்.
  \begin{enumerate}
  \item $M$ என்பது $y$ என்ற ஒரு மாறியாகவும்
    $y \not\ident x$ ஆகவும் இருந்தால், $\Subst{y}{N}{x}$
    என்பது $y$ தான்; ஆகவே அது ரெடெக்ஸ் அல்ல.
  \item $M$ ஆனது $\andi{N_1}{N_2}$, அல்லது
    $\lambd[x][N]$, அல்லது $\ori{i}{!A}{N}$ என்ற
    வடிவில் இருந்தால்,
    $\Subst{M}{\typeof{N}{!A}}{\typeof{x}{!A}}$ உம்
    அதே வெளிப்புற வடிவில் இருக்கும்; ஆகவே அது ரெடெக்ஸ் அல்ல.
  \item $M$ ஆனது $\ande{i}{P}$ என்ற வடிவில் இருந்தால்,
    இரண்டு நேர்வுகளைப் பார்க்கிறோம்.
    \begin{enumerate}
      \item $P$ ஆனது $\andi{P_1}{P_2}$ என்ற வடிவில்
        இருந்தால், $M \ident \ande{i}{\andi{P_1}{P_2}}$
        ஒரு ரெடெக்ஸ்; மேலும் தெளிவாக,
        \[
        \Subst{M}{N}{x} \ident
        \ande{i}{\andi{\Subst{P_1}{N}{x}}{\Subst{P_2}{N}{x}}}
        \]
        என்பதும் ரெடெக்ஸ். இவற்றின் வெட்டு நிலைகள் சமம்.
      \item $P$ ஒற்றை மாறியாக இருந்தால், மாற்றீட்டின்
        விளைவு ரெடெக்ஸ் ஆவதற்கு அது $x$ ஆகவே வேண்டும்;
        $N$ ஆனது $\andi{P_1}{P_2}$ என்ற வடிவில்
        இருக்க வேண்டும். இப்போது
        \[
        \Subst{M}{N}{x} \ident \Subst{\ande{i}{x}}{\andi{P_1}{P_2}}{x},
        \]
        என்பதைப் பார்ப்போம். இது
        $\ande{i}{\andi{P_1}{P_2}}$ ஆகும். அதன்
        வெட்டு நிலை $\len{!A}$.
    \end{enumerate}
  \end{enumerate}
  $\ore{N}{x_1}{N_1}{x_2}{N_2}$ மற்றும் $P Q$
  நேர்வுகள் இதே முறையில் அமைகின்றன.
\end{proof}

% SEGMENT OLP-0687-S04
\begin{lem}
  $M$ ஆனது $M'$ ஆகச் சுருங்குகிறது என்றும்,
  $M$ இன் முறையான ரெடெக்ஸ் உட்சொல் ஒவ்வொன்றுக்கும்
  $\cutrank{M} > \cutrank{N}$ என்றும் வைத்தால்,
  $\cutrank{M} > \maxrank{M'}$.
\end{lem}

\begin{proof}
  நேர்வுகளாகப் பிரித்து நிறுவுகிறோம்.
  \begin{enumerate}
  \item $M$ ஆனது $\ande{i}{\andi{M_1}{M_2}}$
    என்ற வடிவில் இருந்தால், $M'$ என்பது $M_i$.
    $M_i$ இன் ஒவ்வொரு உட்சொல்லும் $M$ இன் முறையான
    உட்சொல்லாகவும் இருப்பதால், கூற்று அமைகிறது.
  \item $M$ ஆனது
    $(\lambd[\typeof{x}{!A}][N])\typeof{Q}{!A}$
    என்ற வடிவில் இருந்தால், $M'$ என்பது
    $\Subst{N}{\typeof{Q}{!A}}{\typeof{x}{!A}}$.
    $M'$ இல் உள்ள ஒரு ரெடெக்ஸைப் பார்ப்போம்.
    அதற்குச் சமமான வெட்டு நிலையுடன் $N$ இல் இருந்தே
    தொடர்புடைய ரெடெக்ஸ் ஒன்று வந்திருக்கும்; அதன்
    நிலை கருதுகோளின்படி $\cutrank{M}$ ஐவிடக் குறைவு.
    அல்லது அதன் வெட்டு நிலை $\len{!A}$ ஆகும்;
    வரையறைப்படி அது
    $\cutrank{(\lambd[x^{!A}][N])Q}$ ஐவிடக் குறைவு.
  \item $M$ ஆனது
    \[
    \ore{\ori{i}{}{\typeof{N}{!A_i}}}
        {\typeof{x_1}{!A_1}}{\typeof{N_1}{!C}}
        {\typeof{x_2}{!A_2}}{\typeof{N_2}{!C}},
    \]
    என்ற வடிவில் இருந்தால்,
    $M' \ident \Subst{N_i}{N}{\typeof{x_i}{!A_i}}$.
    $M'$ இல் உள்ள ஒரு ரெடெக்ஸுக்கு, $N_i$ இலிருந்து
    வந்த அதே வெட்டு நிலை இருக்கலாம்; அது கருதுகோளின்படி
    $\cutrank{M}$ ஐவிடக் குறைவு. இல்லையெனில் அதன்
    வெட்டு நிலை $\len{!A_i}$; வரையறைப்படி அது
    $\cutrank{\ore{\ori{i}{}{\typeof{N}{!A_i}}}
      {\typeof{x_1}{!A_1}}{\typeof{N_1}{!C}}
      {\typeof{x_2}{!A_2}}{\typeof{N_2}{!C}}}$
    ஐவிடக் குறைவு.
  \end{enumerate}
\end{proof}

% SEGMENT OLP-0687-S05
\begin{thm}
  எல்லா நிறுவல் சொற்களையும் இயல்புவடிவமாகக் குறைக்கலாம்;
  எல்லா !!{derivation}s ஐயும் இயல்பு வடிவ
  !!{derivation}s ஆகக் குறைக்கலாம்.
\end{thm}

\begin{proof}
  இரண்டாம் கூற்று முதலாவதிலிருந்து பின்வருகிறது.
  நிறுவல் சொல் $M$ க்கு $m=\maxrank{M}$ எனக் கொண்டு,
  அந்த உச்ச வெட்டு நிலைமீதான முழுத் தொகுத்தறிதலால் முதல் கூற்றை
  நிறுவுகிறோம்.
  \begin{enumerate}

  \item $m=0$ என்றால், $M$ ஏற்கெனவே இயல்புவடிவத்தில் உள்ளது.
  \item மற்றபடி, $M$ இல் வெட்டு நிலை~$m$ உடைய
    ரெடெக்ஸ்களின் எண்ணிக்கை~$n$ மீதான தொகுத்தறிதலைப்
    பயன்படுத்துகிறோம்.
    \begin{enumerate}
    \item $n=1$ என்றால், $m = \cutrank{N} >
      \cutrank{P}$ என்ற நிபந்தனையை நிறைவேற்றும்
      ரெடெக்ஸ்~$N$ ஐத் தேர்கிறோம்; இங்கு $P$ என்பது
      $N$ இன் முறையான ரெடெக்ஸ் உட்சொல். ஒரு சொல்லுக்கு
      முடிவுள்ள எண்ணிக்கையிலேயே உட்சொற்கள் உள்ளதால்,
      அத்தகைய தேர்வு கிடைக்கும்.

      $N'$ என்பது $N$ இன் சுருக்கவிளைவு ஆகட்டும்.
      முந்திய துணைத்தேற்றத்தின்படி
      $\maxrank{N'} < \maxrank{N}$; ஆகவே
      $\cutrank=m$ உடைய ரெடெக்ஸ்களின் எண்ணிக்கையான
      $n$ குறைகிறது. இங்கு அது ஒன்றே இருந்ததால்,
      $m$ உம் குறைந்தது $1$ அளவு குறைகிறது.
      எனவே $m$ க்கான தொகுத்தறிதல் கருதுகோளைப்
      பயன்படுத்தலாம்.
    \item தொகுத்தறிதல் படியில் $n>1$ என்று கொள்வோம்.
      அதே உச்ச வெட்டு நிலை உடைய ரெடெக்ஸ்களில், அதே
      நிலை உடைய முறையான ரெடெக்ஸ் உட்சொல் இல்லாத
      ஒன்றைத் தேர்ந்து சுருக்குகிறோம். முறையும் இதுவே;
      ஆனால் இப்போது $n$ ஒரு நேர்ம எண்ணாகவே குறையலாம்,
      எனவே $m$ மாறாமல் இருக்கலாம். அப்போது $n$
      க்கான தொகுத்தறிதல் கருதுகோளைப் பயன்படுத்துகிறோம்.
    \end{enumerate}
  \end{enumerate}
\end{proof}

இந்த அளவீட்டு வாதம் மேலே பட்டியலிட்ட அறிமுகம்--நீக்கம்
ரெடெக்ஸ்களின் குறைத்தலைக் கையாள்கிறது. முந்திய
குறைப்புப் பிரிவில் விவரிக்கப்பட்ட வரிசைமாற்று
உருமாற்றங்களின் முடிவுறுதலுக்கு இங்கு தனி வாதம் தரப்படவில்லை.

% SEGMENT OLP-0687-S06
வகையிட்ட $\lambd$-சொற்கள் இயல்பாக்கப்படுகின்றன
என்பதன் பொருள்: அத்தகைய ஒவ்வொரு சொல்லுக்கும்
\emph{ஓர் இயல்புவடிவம் இருக்கிறது}. $\lambd$-சொற்களை
$\beta$-குறைத்தல் மூலம் இயல்புவடிவத்துக்குக் கொண்டு
செல்வதைக் கணக்கீடாகவும், இயல்புவடிவத்தை அதன்
மதிப்பாகவும் கொண்டால், வகையிட்ட $\lambd$-கலனத்தில்
ஒவ்வொரு கணக்கீட்டையும் நிறுத்தி மதிப்பு பெறும் ஒரு
குறைத்தல் வழி உண்டு. மேலும், வகை காக்கப்படுவதால்
அம்மதிப்புக்குச் சரியான வகை உண்டு. எடுத்துக்காட்டாக,
$N$ இன் வகை $!A \lif !B$ என்றால், அது வகை~$!A$
உள்ள உள்ளீடுகளில் செயல்பட்டு வகை~$!B$ உள்ள மதிப்புகளைத்
தர வேண்டிய சார்பு. வகை~$!A$ உடைய சொல்~$M$ ஐ
உள்ளீடாக அதன்மீது செயல்படுத்தினால், $(NM)$ ஆனது
வெளியீடான $N'$ என்ற சொல்லாகக் குறைகிறது; அந்த
வெளியீட்டின் வகை~$!B$ என்பது உறுதி.

மேலும் வலுவான பண்பு ஒன்றை மூல உரை இங்கு கூறுகிறது:
வகையிட்ட $\lambd$-கலனச் சொற்கள் ஒரு குறிப்பிட்ட
$\beta$-குறைத்தல் வரிசையில் மட்டும் இயல்புவடிவம்
அடைவதல்ல; உட்சொற்களுக்கு $\beta$-குறைத்தல்களை
\emph{எந்த} வரிசையில் பயன்படுத்தினாலும், இறுதியில்
இயல்புவடிவத்தில் முடியும். இதனை \emph{வலுவான
இயல்பாக்கம்} என்போம். ஆகவே ஒரு வழியில் கணக்கீடு
நின்று மதிப்பைத் தருகிறது என்பதைக் காட்டிலும் வலுவாக,
\emph{எந்த} இயக்க வழியும் இறுதியில் மதிப்பைத் தருகிறது.
இந்த வலுவான பண்புக்கு மேலுள்ள ஒருவழி இயல்பாக்க
வாதமே தனியாக நிறுவலாகாது.

% SEGMENT OLP-0687-S07
வெவ்வேறு வழிகளில் கணக்கிடும்போது \emph{ஒரே}
மதிப்புதான் கிடைக்கிறது என்பதை எப்படி அறிவது?
இதுவரை வகை காக்கப்படுவதால் எல்லா மதிப்புகளும்
ஒரே வகையுடையவை என்பதை மட்டுமே அறிந்தோம்.
வகையிட்ட $\lambd$-கலனத்துக்கு மற்றொரு முக்கியப்
பண்பு உள்ளது: அது \emph{சங்கமப் பண்பு} அல்லது
\emph{சர்ச்--ரோசர் பண்பு} உடையது. $N \red N_1$
மற்றும் $N \red N_2$ என்றால், $N_1 \red N'$ உம்
$N_2 \red N'$ உம் ஆகும் வகையில் ஒரு சொல்~$N'$
உண்டு. $N \red N'$ என்பதன் பொருள்: பூச்சியம் அல்லது
அதற்கு மேற்பட்ட $\redone$-குறைத்தல் படிகளால் $N$
இலிருந்து $N'$ கிடைக்கிறது. $N_1$ இயல்புவடிவத்தில்
இருந்தால், $N_1 \red N'$ ஆகும் ஒரே சொல்லான $N'$
என்பது $N_1$ தான்;
அது பூச்சிய $\redone$-படிகளால் கிடைக்கிறது.
ஆகவே $N_1$ உம் $N_2$ உம் இயல்புவடிவத்தில் இருந்தால்,
$N_1 = N' = N_2$. வேறு சொற்களில், $\red$
வலுவாக இயல்பாக்கப்படுமானால், எல்லாக் குறைத்தல்
வழிகளும் இறுதியில் இயல்புவடிவத்தை அடைகின்றன.
$\red$ சர்ச்--ரோசர் பண்பும் உடையதானால்,
இரு இயல்புவடிவங்களும் சமம். ஆகவே இவ்விரு பண்புகளும்
இருக்கும் நிலையில் வகையிட்ட $\lambd$-கலனக்
கணக்கீடு எப்போதும் நிற்கும்; கணக்கீட்டு வழி எதுவானாலும்
அதன் மதிப்பு, அதாவது இயல்புவடிவம், ஒரேதுதான்.

% SEGMENT OLP-0687-S08
ஒரு தொடர்புக்குச் சங்கமப் பண்பு இருப்பதை நிறுவுவது
பொதுவாகக் கடினம். ஆனால் பின்வரும் பலவீனமான
நிபந்தனையை நிறுவலாம்:

\begin{prop}[பலவீனச் சர்ச்--ரோசர் பண்பு]
$N \redone N_1$ மற்றும் $N \redone N_2$ என்றால்,
$N_1 \red N'$ உம் $N_2 \red N'$ உம் ஆகும் வகையில்
ஒரு சொல் $N'$ உண்டு.
\end{prop}

\begin{proof}
  $N$ இல் உள்ள ரெடெக்ஸ் $M_1$ அல்லது $M_2$ ஐ
  முறையே அதன் சுருக்கவிளைவால் மாற்றுவதால் $N_1$,
  $N_2$ கிடைக்கின்றன. இரண்டிலும் அதே ரெடெக்ஸ்
  நிகழ்வு தேர்ந்தெடுக்கப்பட்டால் இரு விளைவுகளும் ஒன்றே.
  $M_1$ உம் $M_2$ உம் $N$ இன் வேறு உட்சொல்
  நிகழ்வுகளாக இருந்தால்,
  ஒன்று மற்றொன்றின் உட்சொல்லாக இருக்கும்; அல்லது
  அவை வெவ்வேறு, ஒட்டாத இடங்களில் இருக்கும்.
  பின்னைய நேர்வைப் பார்ப்போம்: $N$ ஆனது
  $N(M_1,M_2)$ என்ற வடிவில் உள்ளது. $M_1$ ஐ அதன்
  சுருக்கவிளைவு $M_1'$ ஆல் மாற்றினால்
  $N(M_1',M_2)$ ஆகிய $N_1$ கிடைக்கும்;
  $M_2$ ஐ $M_2'$ ஆல் மாற்றினால் $N(M_1,M_2')$
  ஆகிய $N_2$ கிடைக்கும். இரண்டும்
  $N(M_1',M_2')$ ஆகக் குறைகின்றன.

  மற்ற நேர்வில் $M_2$, $M_1$ இன் உட்சொல் என்று
  கொள்வோம்; எதிர்த்திசை நேர்வு சமச்சீரானது. $M_1$
  எந்த வகை ரெடெக்ஸ் என்பதைப் பொறுத்து நேர்வுகள்
  வேறுபடும். எடுத்துக்காட்டாக,
  $M_1 \ident \proj{1}{\pair{O_1}{O_2}}$
  என்றால் அது $O_1$ ஆகக் குறையும். $M_2$ ஆனது
  $O_1$ இன் உட்சொல் என்றால், அதைக் குறைத்த பின்
  $O_1'$ கிடைக்கும். முதலில் $M_1$ ஐயும் பின்னர்
  $M_2$ ஐயும் குறைத்தாலும் $O_1'$ கிடைக்கும்.
  முதலில் $M_2$ ஐக் குறைத்தால்
  $\proj{1}{\pair{O_1'}{O_2}}$ கிடைக்கும்;
  இதுவும்~$O_1'$ ஆகக் குறையும். $M_2$ கைவிடப்படும்
  இரண்டாம் கூறில் இருந்தால், எந்த வரிசையிலும்
  இறுதியில் தேர்ந்தெடுக்கப்பட்ட முதல் கூறே கிடைக்கும்.
  சார்பு-செயல்படுத்தல்,
  கூட்டல்-வகை நேர்வுகளிலும் உள்ளேயுள்ள குறைத்தலை
  மாற்றீட்டுக்குப் பிறகு அதன் நகல்களில் தொடர்ந்து
  செய்ய வேண்டும். இந்நேர்வுகளின் முழு விளக்கக்
  கணக்கை உறைந்த மூலம் இங்கு தரவில்லை.
\end{proof}

% SEGMENT OLP-0687-S09
\begin{lem}[நியூமனின் துணைத்தேற்றம்]
$\red$ வலுவாக இயல்பாக்கப்படுவதுடன் பலவீனச்
சர்ச்--ரோசர் பண்பையும் கொண்டிருந்தால், அது சங்கமப்
பண்புடையது.
\end{lem}

\begin{proof}
  $\red$ வலுவாக இயல்பாக்கப்படுவதால், எந்தக்
  குறைத்தல் தொடரும் இயல்புவடிவத்தில் முடியும்.
  இயல்புவடிவங்கள் தனித்தனியாக ஒரே முடிவைக்
  கொடுப்பதும் $\red$ சங்கமப் பண்புடையதாக
  இருப்பதும் இந்நிலையில் சமமானவை. ஆகவே ஒரு சொல்
  $N$ க்கு $N'$ மற்றும் $N''$ என்ற இரு வேறு
  இயல்புவடிவங்கள் உள்ளன எனக் கொள்வோம். அவை
  பின்வரும் குறைத்தல் தொடர்களின் விளைவுகள்:
  \begin{align*}
  N & = N_1' \redone N_2' \redone \dots \redone N_k' = N' \text{ மற்றும்}\\
  N & = N_1'' \redone N_2'' \redone \dots \redone N_l'' = N''
  \end{align*}
  இத்தொடர்களின் நீளம் $>0$ ஆகவே வேண்டும்; இல்லை
  என்றால் $N$ ஏற்கெனவே இயல்புவடிவத்தில் இருக்கும்.

  $N_1' = N_1''$ என்றால், அந்தச் சொல்லுக்கே
  $N'$ உம் $N''$ உம் என்ற இரு வேறு இயல்புவடிவங்கள்
  உள்ளன. $N_1' \neq N_1''$ என்றால், பலவீனச்
  சர்ச்--ரோசர் பண்பால் $N_1' \red N'''$ உம்
  $N_1'' \red N'''$ உம் ஆகும் ஓர் $N'''$ உண்டு.
  $N' \neq N''$ என்பதால், $N''' \neq N'$
  அல்லது $N''' \neq N''$ என்பது உறுதி. இங்கு
  வலுவான இயல்பாக்கத்தால் அந்தப் பொதுப் பின்விளைவுக்கும்
  ஓர் இயல்புவடிவம் உண்டு; அது முன்னர் கிடைத்த இரு
  வேறு இயல்புவடிவங்களில் குறைந்தது ஒன்றிலிருந்து
  வேறுபடும். எனவே $N_1'$ அல்லது $N_1''$ இல்
  ஒன்றுக்கு இரு வேறு இயல்புவடிவங்கள் உள்ளன.
  அதாவது, $N$ க்கு இரு வேறு இயல்புவடிவங்கள்
  இருந்தால், $N \redone N_1$ ஆகவும் $N_1$
  க்கும் இரு வேறு இயல்புவடிவங்கள் இருக்கவும்
  ஓர் $N_1$ கிடைக்கும். இதேபோல் தொடர்ந்தால்
  $N \redone N_1 \redone N_2 \redone \dots$
  என்ற முடிவிலாத குறைத்தல் தொடர் கிடைக்கும்;
  ஆனால் $\redone$ வலுவாக இயல்பாக்கப்படுவதால்
  அது இயலாது.
\end{proof}

\end{document}
